\documentclass[10pt,a4paper]{amsart}
\usepackage[margin=19mm]{geometry}
\usepackage{amsmath,amssymb,mathtools}
\usepackage[scr=rsfs,cal=euler]{mathalfa}
\usepackage{xfrac}
\usepackage[unicode,hidelinks,bookmarks=false]{hyperref}
\newcommand{\ie}{i.e.\ }
\newcommand{\cf}{cf.\ }
\newcommand{\resp}{respectively\ }
\newcommand{\Subsection}{\subsection}
\providecommand{\nobreakdash}{\nobreak\hbox{-}}
\setlength{\emergencystretch}{2em}
\multlinegap0pt
\usepackage{bm}
\usepackage{bbm}
\usepackage{dsfont}
\usepackage{xspace}
\usepackage{cancel}
\usepackage{mathtools}
\usepackage{amsthm}
\makeatletter
\@ifundefined{lemma}{\newtheorem{lemma}{Lemma}}{}
\makeatother
\def\R{\mathbb{R}}
\def\nn{\nonumber}
\title[A gradient flow model for the Gross--Pitaevskii problem: Mat]{A gradient flow model for the Gross--Pitaevskii problem: Mathematical and numerical analysis}
\author{Tianyang Chu, Xiaoying Dai, Jing Wu, Aihui Zhou}
\date{Source: arXiv:2510.27235v1 (CC BY 4.0)}
\begin{document}
\maketitle

\noindent\emph{Source and licence.} A gradient flow model for the Gross--Pitaevskii problem: Mathematical and numerical analysis. Version arXiv:2510.27235v1 (CC BY 4.0), distributed under the Creative Commons Attribution 4.0 International licence, as recorded in the arXiv licence metadata for this exact version and shown on the abstract page https://arxiv.org/abs/2510.27235v1. Full rights notice: licenses/arxiv-2510.27235-CC-BY-4.0.txt.\par\smallskip

\noindent\textit{Editorial scope.} Counted unit: the Lemma whose source label is \texttt{lem: key lemma} in the article (source0.tex lines 479--488, source label \texttt{lem: key lemma}), delivered together with the article's own complete original proof of it (source0.tex lines 489--569, one \texttt{proof} environment of 81 lines). No mathematics is written, repaired or re-derived: statement and proof are the article's own text line for line. Required context retained uncounted: none. Every definition, display and label of those lines is retained unchanged. Omitted from this package and not redistributed: 1--478, 570--1493. Nothing inside a retained excerpt is omitted or altered: every retained body is delivered exactly as the article's lines are written, so no omission receipt of any kind accompanies this package. Citations: the retained excerpts carry no citation command at all, so no bibliography entry of the article is transcribed and no reference is redistributed. Reference adaptation: none was needed and none was made; every reference inside the delivered unit resolves inside it, so no unresolved reference or citation is produced. Printed slips kept literal: no printed word, symbol, hypothesis or number of the counted statement or of its proof was altered. Layout and class: the article's own class is \texttt{siamart250211} with packages including microtype, amssymb, bm, epstopdf, xcolor, tikz, float, graphicx, subfig, algorithmic, amsfonts, mathtools, nicefrac, mathrsfs; the wrapper is the lane's pinned \texttt{amsart} editorial wrapper at 10pt on A4, which restates the theorem-like environments the retained text needs and adds the article's own macro definitions that the retained text needs (\texttt{\textbackslash R}, \texttt{\textbackslash nn}), with extra wrapper packages (bm, bbm, dsfont, xspace, cancel, mathtools). The delivered numbering is the unit's own. Assets: no retained span contains a figure environment, a graphics-inclusion command, a PDF-page inclusion, a table or a TikZ picture; no image file is copied into this package, the package declares no asset dependency and no image file is redistributed with it. Licence: version arXiv:2510.27235v1 of this article is distributed under the Creative Commons Attribution 4.0 International licence, as recorded in the arXiv licence metadata for this exact version and shown on the abstract page https://arxiv.org/abs/2510.27235v1. A full rights notice is delivered at licenses/arxiv-2510.27235-CC-BY-4.0.txt. Characters: every retained excerpt is pure ASCII, so the delivered engine, XeLaTeX with its default Latin Modern text font, reports no missing character.
\begin{lemma}\label{lem: key lemma}
Denote $g(t) = \frac12\|\phi_t(\cdot,t) \|_{L^2(\Omega)}^2$, then for any $\epsilon \in (0, \lambda_2 - \lambda_{\mathrm{GS}})$, there exists $T_{\epsilon} > 0$ such that 
	\begin{equation}\label{equ: key inequality}
		g'(t) \le - 2(\lambda_2 - \lambda_{\text{GS}} - \epsilon) g(t) \qquad t \ge T_{\epsilon},
	\end{equation}
	where $\lambda_2$ is the second eigenvalue of the following linear eigenvalue problem: 
	\begin{equation*}
		- \Delta u + V u + \beta |\phi_{\text{GS}}|^2 u = \lambda u.
	\end{equation*}
\end{lemma}

\begin{proof}
	Note that 
\begin{equation*}
	\phi_{tt} = \Delta \phi_t - V \phi_t - 3 \beta \phi^2 \phi_t + \left( \frac{\text{d}}{\text{d}t} \mu[\phi] \right) \phi + \mu[\phi]\phi_t,
\end{equation*}
where $\partial_{tt}\phi$ is understood in $H^{-1}$ sense. We have that 
\begin{align}
	g'(t) & = \langle\phi_t, \phi_{tt}\rangle_{H^{-1}(\Omega) \times H_0^1(\Omega)} \nn \\
	& = -(\nabla \phi_t,  \nabla \phi_t)_{L^2(\Omega)} - (V\phi_t,\phi_t)_{L^2(\Omega)}  - \beta (\phi^2 \phi_t, \phi_t)_{L^2(\Omega)} \underbrace{- 2\beta (\phi^2 \phi_t, \phi_t)_{L^2(\Omega)}}_{\le 0} \nn \\
	& \quad + \left( \frac{\text{d}}{\text{d}t} \mu[\phi] \right) \underbrace{(\phi, \phi_t)_{L^2(\Omega)}}_{ = 0} + \mu[\phi] (\phi_t,\phi_t)_{L^2(\Omega)} \nn \\
	& \le - \int_{\Omega} \Big( |\nabla \phi_t|^2 + V|\phi_t|^2 + \beta|\phi|^2 |\phi_t|^2 \Big) \ \text{d}x + \mu[\phi]\|\phi_t\|_{L^2(\Omega)}^2.
\end{align}

Denote 
\begin{equation*}
a_{\phi}(u,v):= \int_{\Omega} \Big( \nabla u \cdot \nabla v + V uv + \beta|\phi|^2 uv \Big) \ \text{d}x,
\end{equation*}
\begin{equation*}
	C_{\text{inf}}:= \liminf_{t \to \infty} \frac{a_{\phi}(\phi_t,\phi_t)}{\|\phi_t\|_{L^2(\Omega)}^2}.
\end{equation*}
Assume that $(t_n)_{n \in \mathbb{N}}$ with $t_n \to \infty$ satisfies
\begin{equation*}
	\lim_{n \to \infty} \frac{a_{\phi(t_n)}(\phi_t(t_n),\phi_t(t_n))}{\|\phi_t(t_n)\|_{L^2(\Omega)}^2} = C_{\text{inf}}.
\end{equation*}
Let
\begin{equation*}
	z_n:= \frac{\phi_t(t_n)}{\|\phi_t(t_n)\|_{L^2(\Omega)}},
\end{equation*}
we see that $\{z_n\}_{n \in \mathbb{N}}$ is a bounded sequence in $H_0^1(\Omega)$ and
\begin{align*}
	&\left| \int_{\Omega} |\phi(t_n)|^2 z_n^2 \ \text{d}x - \int_{\Omega} |\phi_{\text{GS}}|^2 z_n^2 \ \text{d}x    \right| \\
	 & \le \| \phi(t_n) + \phi_{\text{GS}}\|_{L^4(\Omega)}\| \phi(t_n) - \phi_{\text{GS}}\|_{L^4(\Omega)}\|z_n^2\|_{L^2(\Omega)}  \to 0.
\end{align*}
Hence, 
\begin{equation*}
	\lim_{n \to \infty} \underbrace{ \int_{\Omega} \Big( |\nabla  z_n|^2 + V| z_n|^2 + \beta|\phi_{\text{GS}}|^2 |z_n|^2 \Big) \ \text{d}x}_{=: a_{\phi_{\text{GS}}}(z_n,z_n)} = C_{\text{inf}}.
\end{equation*}

Since $\{z_n\}_{n \in \mathbb{N}}$ is bounded in $H^1(\Omega)$, there exists a weak limit 
\begin{equation*}
	z_n \rightharpoonup \hat{z} \qquad \text{weakly in $H^1(\Omega)$}.
\end{equation*}
This indicates $\|\hat{z}\|_{L^2(\Omega)} = 1$ and
\begin{align*}
	&\left| (\hat{z}, \phi_{\text{GS}})_{L^2(\Omega)} - (z_n, \phi(t_n))_{L^2(\Omega)}  \right| \\
	 &\le \|\hat{z} - z_n\|_{L^2(\Omega)}\|\phi_{\text{GS}}\|_{L^2(\Omega)} + \|z_n\|_{L^2(\Omega)}\|\phi_{\text{GS}} - \phi(t_n)\|_{L^2(\Omega)}  \to 0,
\end{align*} 
which leads to
\begin{equation*}
	(\hat{z}, \phi_{\text{GS}})_{L^2(\Omega)} = 0.
\end{equation*}
Lower semicontinuity of weakly converging sequence implies 
\begin{equation*}
	C_{\text{inf}} = \lim_{n \to \infty}a_{\phi_{\text{GS}}}(z_n,z_n) \ge a_{\phi_{\text{GS}}}(\hat{z}, \hat{z}) \ge \inf_{v \in \text{span}\{\phi_{\text{GS}}\}^\bot} \frac{a_{\phi_{\text{GS}}}(v, v)}{\|v\|_{L^2(\Omega)}^2}.
\end{equation*}

With the Courant-Fischer theorem we get
\begin{equation*}
	C_{\text{inf}} \ge \lambda_2,
\end{equation*}
where $\lambda_2$ is the second eigenvalue of linear eigenvalue problem: Find $(u,\lambda) \in H_0^1(\Omega) \times \R$ such that
\begin{equation*}
	- \Delta u + V u + \beta |\phi_{\text{GS}}|^2 u = \lambda u.
\end{equation*}
	
With the fact $\mu[\phi(t)] \to \lambda_{\text{GS}}$ as $t \to \infty$, there exists $T_{\epsilon} > 0$ such that for $t \ge T_{\epsilon}$
\begin{equation}\label{equ: 1}
	\mu[\phi(t)] \le \lambda_{\text{GS}} + \frac{\epsilon}{2}
\end{equation}
and 
\begin{equation}\label{equ: 2}
	\frac{a_{\phi}(\phi_t,\phi_t)}{\|\phi_t\|_{L^2(\Omega)}^2} \ge \lambda_2 - \frac{\epsilon}{2}.
\end{equation}

Combining \eqref{equ: 1} and \eqref{equ: 2}, we have that for $t \ge T_{\epsilon}$, 
\begin{align*}
	g'(t) & \le - \int_{\Omega} \Big( |\nabla \phi_t|^2 + V|\phi_t|^2 + \beta|\phi|^2 |\phi_t|^2 \Big) \ \text{d}x + \mu[\phi]\|\phi_t\|_{L^2(\Omega)}^2  = \|\phi_t\|_{L^2(\Omega)}^2 \left( \mu[\phi] - \frac{a_{\phi}(\phi_t,\phi_t)}{\|\phi_t\|_{L^2(\Omega)}^2} \right) \\
	& \le - (\lambda_2 - \lambda_{\text{GS}} - \epsilon) \|\phi_t\|_{L^2(\Omega)}^2 = -2(\lambda_2 - \lambda_{\text{GS}} - \epsilon)g(t).
\end{align*}
This completes the proof.
\end{proof}

\end{document}
